\documentclass[a4paper,11pt]{ltjsarticle}
\usepackage[haranoaji]{luatexja-preset}
\usepackage[margin=25mm]{geometry}
\usepackage{fix-cm}
\usepackage{amsmath,amssymb,amsthm,mathtools,booktabs,array}
\usepackage[unicode,hidelinks]{hyperref}
\usepackage{xurl}
\theoremstyle{definition}
\newtheorem{theorem}{定理}[section]
\newtheorem{proposition}[theorem]{命題}
\newtheorem{lemma}[theorem]{補題}
\newtheorem{corollary}[theorem]{系}
\theoremstyle{definition}\newtheorem{definition}[theorem]{定義}
\theoremstyle{remark}\newtheorem{remark}[theorem]{注意}
\newcommand{\PP}{\mathbf P}\newcommand{\CC}{\mathbf C}
\newcommand{\cO}{\mathcal O}\newcommand{\cP}{\mathcal P}
\newcommand{\Aut}{\operatorname{Aut}}\newcommand{\PGL}{\operatorname{PGL}}
\newcommand{\Stab}{\operatorname{Stab}}\newcommand{\rank}{\operatorname{rank}}
\newcommand{\diver}{\operatorname{div}}\newcommand{\Pf}{\operatorname{Pf}}
\title{射影四空間のPoisson分類の構造還元\\\large 標識因子・modular復元・相対変形}
\author{}
\date{2026年11月9日\quad v4・二次因子と磁気留数}
\begin{document}\maketitle
\begin{abstract}
射影空間上の一般に非退化なPoisson構造について、退化因子の既約成分が生成する閉対数一形式を用い、
$r-1$次元の標準的な可換対称性代数を構成する。退化因子の被約性や正規交叉性は仮定しない。
射影四空間では、射影線形軌道の次元は$25-r$以下となる。
また四次元では、退化因子を固定したPoisson構造の分類を、閉有理二形式の線形条件と
Pfaffianの二次条件へ置き換える。これは単なる必要条件でなく、一般階数四の全点についての同値である。
さらに全階数四を三つの入口へ分け、超平面を持つ全群を一意なアフィン二次模型で統合する。
Pfaffianとmodular場を固定した線形・二次復元、共通の共鳴層別、因子変形の一次・二次持上げ、
Pfaffianグラフの閉包、および階数二のfoliationと零因子による入口を与える。
二次＋三次ではcanonical cofactorと全31場を同じ斉次復元へ接続する。
被約平面五次の全sourceと磁気shear核を因子留数で記述し、因子数$\nu$に対し
磁気商の次元$\nu-1$、その零境界と軌道次元を導く。
これらを全支持の認識と成分の極大性を分離する基盤とする。
本稿は147成分という個数の新証明でも、その全証明の独立検証でもない。
以下の基本原理には通常の証明を与えるが、それ自体の文献上の新規性は未確認であり、優先権を主張しない。
\end{abstract}

\section{列挙ではなく構造を中心に置く}
一般階数四のPoisson構造$\pi$を持つ$X=\PP^4$に対し、
\[
 D_\pi=\operatorname{div}(\pi\wedge\pi/2)\in|-K_X|=|\cO_X(5)|
\]
を\emph{重複度を含む}実退化因子とする。
$U=X\setminus D_\pi$上の逆二形式$\omega=\pi^{-1}$は閉である。
この対応と退化因子・対数形式の基本理論は\cite[\S5]{Pym}に説明されている。

本稿の整理の軸は
\[
 \pi\quad\longmapsto\quad
 \bigl(D_\pi,\ \mathfrak a_\pi,\ \omega\bigr)
\]
である。$D_\pi$は極の幾何、$\mathfrak a_\pi$は因子から強制される射影線形対称性、
$\omega$はその幾何の上に残る閉形式の情報を表す。
三者が互いに独立な自由データであるという意味ではない。

\section{退化因子の因子から可換対称性を作る}
以下では$X=\PP^{2n}$、$n\ge1$とする。
一般に非退化なPoisson構造の退化因子を
\[
 D=\sum_{i=1}^{r}m_iD_i,\qquad D_i=V(f_i),\qquad
 d_i=\deg f_i,\qquad \sum_i m_id_i=2n+1
\]
と書く。$f_i$は互いに相異なる既約斉次式、$m_i>0$である。

\begin{lemma}\label{lem:tangent}
各$D_i$の被約イデアルはPoissonイデアルである。
\end{lemma}
\begin{proof}
局所座標で$\pi^n/n!=F\partial_1\wedge\cdots\wedge\partial_{2n}$と書く。
任意の正則関数$g$のHamiltonian場$H_g=\pi^\sharp(dg)$は$\pi$を保存する。
従って$H_g(F)=F\diver(H_g)$であり、$(F)$はHamiltonian導分に関して安定である。
滑らかな多様体の局所環は正則局所環なので一意分解環である。
$F=u\prod_i g_i^{m_i}$を局所的に因数分解し、$H_g(F)/F$の各素因子上の極を比較すると
$g_i\mid H_g(g_i)$を得る。標数零により$m_i\ne0$が使える。
よって各既約因子のイデアルが安定である。
\end{proof}

\begin{definition}
$D_{\mathrm{red}}$の次数関係から作るベクトル空間を
\[
 W_D=\left\{a=(a_1,\ldots,a_r)\in\CC^r\;\middle|\;
              \sum_i a_id_i=0\right\}
\]
とし、$a\in W_D$に対して
\[
 \eta_a=\sum_i a_i\frac{df_i}{f_i}
\]
と置く。これは$X$上の閉有理一形式である。
\end{definition}

斉次座標のEuler場を$E$とすると、$\iota_E\eta_a=\sum_i a_id_i=0$であり、
また$\eta_a$はスカラー作用に不変なので射影空間へ降下する。
より直接には、線形式$\ell\ne0$のチャート上で$g_i=f_i/\ell^{d_i}$と置けば
$\eta_a=\sum_i a_i\,dg_i/g_i$であり、異なるチャートで一致する。
$f_i$を非零定数倍しても変わらない。

\begin{theorem}[因子由来の標準的な可換対称性]\label{thm:abelian}
写像
\[
 W_D\longrightarrow H^0(X,T_X),\qquad a\longmapsto X_a:=\pi^\sharp\eta_a
\]
は単射であり、その像$\mathfrak a_\pi$は$r-1$次元の可換Lie部分代数である。
各$X_a$は$\pi$と各$D_i$を保存する。
さらに
\[
 b_\pi(a,b):=\pi(\eta_a,\eta_b)
\]
は$X$上の定数であり、$W_D$上の交代双線形形式を定める。
この結論には$D$の被約性、正規性、正規交叉性を仮定しない。
\end{theorem}
\begin{proof}
補題\ref{lem:tangent}より$\pi^\sharp(dg_i)/g_i$は正則ベクトル場なので$X_a$は全域で正則である。
$\eta_a$は閉であり、極のない単連結小開集合ではある正則関数の微分になる。
従って$X_a$はそこでHamiltonianで、$\mathcal L_{X_a}\pi=0$。
両辺は正則なので全域で成り立つ。Poisson対称性は退化因子を保存し、その各既約成分も保存する。

局所的に$i\ne j$なら$\{g_i,g_j\}$は$g_i$でも$g_j$でも割れる。
相異なる因子なので積$g_ig_j$で割れ、従って
$\pi(dg_i/g_i,dg_j/g_j)$は正則である。$i=j$の項は零である。
したがって$b_\pi(a,b)$は大域正則関数であり、射影空間の完備性から定数である。
局所Hamiltonianの交換子公式により
\[
 [X_a,X_b]=\pi^\sharp d\bigl(b_\pi(a,b)\bigr)=0
\]
となる。

最後に$X_a=0$とする。$\pi$は$U$上で非退化なので$\eta_a=0$。
$D_i$の一般滑らかな点で他の因子を避けると、その留数は$a_i$である。
従って全$a_i=0$であり、単射である。$\dim W_D=r-1$から結論を得る。
\end{proof}

\begin{remark}
標準的であるのは対$\,(D,\pi)\,$から得るLie部分代数であり、因子$D$だけでその埋込みまで決まるわけではない。
また$r-1$は\emph{大域ベクトル場のベクトル空間としての次元}であり、各点での評価の階数とは限らない。
それらが生成する連結解析群の作用を、追加の確認なしに代数的なトーラスや加法群の作用と同一視しない。
一般階数二のPoisson構造では単射性の証明が失われるため、この定理をそのまま適用しない。
\end{remark}

\begin{corollary}[軌道次元の一様上界]\label{cor:orbit}
$X=\PP^4$について
\[
 \dim\Stab_{\PGL_5}([\pi])\ge r-1,\qquad
 \dim(\PGL_5\cdot[\pi])\le 25-r.
\]
\end{corollary}
\begin{proof}
$H^0(\PP^4,T_{\PP^4})=\mathfrak{pgl}_5$であり、
$\mathfrak a_\pi$は$\pi$そのものを保存するので、その射影類の安定化Lie代数にも入る。
$\dim\PGL_5=24$から従う。
\end{proof}

したがって、二因子型では一つの非零な射影線形対称性が必ず存在し、軌道次元は23以下となる。
三因子型では二次元の可換対称性が存在し、軌道次元は22以下となる。
これで個々のJordan型を調べる\emph{前に}、対称性を用いる分類の理由が分かる。
ただし上界の等号や族の成分性は別の証明を必要とする。

\section{固定した退化因子上では線形問題と二次問題に分離する}
この節では$X$を任意の滑らかな射影四次元多様体、$L=K_X^{-1}$とする。
四次元の外積双対による標準同型を
\[
 \Phi:\wedge^2T_X\xrightarrow{\sim}\Omega_X^2\otimes L
\]
と書く。$\alpha\in H^0(X,\Omega_X^2\otimes L)$に対して
\[
 q(\alpha)=\frac12\alpha\wedge\alpha
       \in H^0(X,K_X\otimes L^2)=H^0(X,L)
\]
と定める。

\begin{lemma}\label{lem:inverse}
$\alpha=\Phi(\pi)$なら$q(\alpha)=\pi\wedge\pi/2$を標準同一視の下で満たす。
$q(\alpha)\ne0$の開集合上で$\pi$の逆形式は$\alpha/q(\alpha)$である。
従って$\pi$が一般に非退化なら
\[
 [\pi,\pi]=0\quad\Longleftrightarrow\quad
 d\left(\frac{\alpha}{q(\alpha)}\right)=0.
\]
\end{lemma}
\begin{proof}
局所体積形式を$\mu=dz_1\wedge\cdots\wedge dz_4$とし、対応する$L$の枠を固定する。
交代形式の標準形で$\pi=a\partial_1\wedge\partial_2+b\partial_3\wedge\partial_4$なら
$\alpha=a\,dz_3\wedge dz_4+b\,dz_1\wedge dz_2$、$q(\alpha)=ab$である。
したがって$\alpha/q(\alpha)=a^{-1}dz_1\wedge dz_2+b^{-1}dz_3\wedge dz_4$となる。
非退化双ベクトルのJacobi恒等式と逆形式の閉性の同値は、逆行列の微分
$d(P^{-1})=-P^{-1}(dP)P^{-1}$を各三添字で反対称化すれば得られる。
非退化開集合で成り立つJacobi恒等式は正則テンソルの等式として全域へ延長する。
\end{proof}

非零の$F\in H^0(X,L)$を固定し、$D=V(F)$とする。局所的に
\[
 d_F\alpha:=F\,d\alpha-dF\wedge\alpha
\]
と定める。$L$の枠を変えて$F,\alpha$が同じ可逆関数の逆数倍になると、
追加微分項は相殺し$d_F\alpha$はその可逆関数の逆数の二乗倍になる。
従ってこれは大域的な\emph{線形}写像
\[
 d_F:H^0(X,\Omega_X^2\otimes L)\longrightarrow
       H^0(X,\Omega_X^3\otimes L^2)
\]
を定める。$V_F:=\ker d_F$と置く。

\begin{theorem}[固定因子の線形・二次表示]\label{thm:fibre}
実退化因子が$D$に等しい一般階数四Poisson構造の射影化された集合は、$\Phi$により
\[
 \left\{[\alpha]\in\PP(V_F)\;\middle|\;
       q(\alpha)\in\CC^{\times}F\right\}                 \tag{3.1}
\]
と一対一に対応する。
これは$\PP(V_F)$の中で、$q(\alpha)$の$H^0(X,L)/\CC F$への像を零とする
斉次二次方程式を課し、$q(\alpha)=0$の閉部分を除いたものである。
\end{theorem}
\begin{proof}
実退化因子が$D$に等しいことは、非零切断$q(\alpha)$が$F$の非零定数倍であることと同値である。
比は零点・極を持たない大域正則関数なので定数である。
$q(\alpha)=cF$なら補題\ref{lem:inverse}によりJacobi恒等式は
$d(\alpha/F)=0$と同値であり、これは$d_F\alpha=0$に等しい。
逆に(3.1)から同じ議論でPoisson構造を復元でき、退化因子も正確に$D$となる。
$q$は二次写像なので最後の表示も従う。
\end{proof}

\begin{remark}
この表示は「閉じた逆形式が得られそうだ」という必要条件ではなく、全点での同値である。
被約因子に限らない。しかし本稿では非被約基底上の変形関手や全Jacobiスキーム構造との同値までは主張しない。
また$D_{\mathrm{red}}$ではなく$F$の\emph{重複度}を保持することが本質的である。
$q(\alpha)=0$を除く条件を落とすとrank二の点が混入し、固定退化因子の分類ではなくなる。
\end{remark}

\section{全階数四を分ける三つの幾何学的入口}\label{sec:global-entry}
\begin{theorem}[次数五による全入口]\label{thm:three-entrances}
一般階数四の$\PP^4$上のPoisson構造の実退化因子$D$は、次のいずれかに属する。
\begin{enumerate}
\item 超平面因子を持つ。重複度は任意であり、$D=5H$を含む。
\item 超平面因子を持たず可約である。この場合は$D=Q+C$で、$Q,C$はそれぞれ既約二次、既約三次の相異なる因子であり、$D$は被約である。
\item $D$は被約既約五次因子である。さらに正規の場合と非正規の場合を分ける。
\end{enumerate}
\end{theorem}
\begin{proof}
$D=\sum m_iD_i$、$d_i=\deg D_i$と書けば$\sum m_id_i=5$である。
$d_i=1$があれば第一の場合。そうでなければ全$d_i\ge2$である。
二つ以上の因子があれば次数の和から、重複度を含めて$2+3$しかない。
一つの因子なら$m_1d_1=5$と$d_1\ge2$から$m_1=1,d_1=5$である。
\end{proof}

これは成分の分割ではなく、\emph{全点を漏らさない分類問題の入口}である。
同じ既約成分は複数の入口の点を含み得るので、各入口で得た族の個数を足してはならない。
とくに被約因子数$r=1$には既約五次だけでなく$5H$もある。
$\mathfrak a_\pi=0$から、他の対称性まで存在しないとは言えない。

% Shared by the classification and structural manuscripts.
\section{標識超平面の共通復元と次数境界}\label{sec:marked-common}
\begin{theorem}[標識超平面の一意なアフィン模型]\label{thm:marked-affine}
$H=V(h)\subset\PP^4$を固定する。$H$をPoisson不変にする射影Poisson双ベクトルと、
$\PP^4\setminus H\simeq\mathbb A^4$上の係数次数二以下のPoisson双ベクトル$P$は自然に一対一に対応する。
対応する斉次二次liftは$h$をCasimirとするものとして一意である。
一般階数四なら、$f=\Pf(P)\ne0$として
\begin{equation}
 D_\pi=V\bigl(h^5f(x/h)\bigr),\qquad
 \operatorname{mult}_H D_\pi=5-\deg f.                 \label{eq:marked-divisor}
\end{equation}
ここで$\deg f\le4$である。したがって単純$H$と重複$H$の入口は、同じ模型で
$\deg f=4$と$\deg f\le3$を分けたものである。
\end{theorem}
\begin{proof}
Euler列の二外積と$H^1(\PP^4,T)=0$により、射影双ベクトルは$\mathbb A^5$上の斉次二次lift $\Pi$を持つ。
liftの差は$E\wedge X$（$E$はEuler場、$X$は線形場）である。
$V=\Pi^\sharp dh$と置く。$H$の不変性は、$h=0$で$V$がEuler方向に属することを意味する。
従って$V_i x_j-V_jx_i=0$が$h$を法として成り立つ。
多項式環の一意分解性から、ある一次式$a$と一次式$W_i$により
$V_i=ax_i+hW_i$（$1\le i\le4$）と書ける。
$X(h)=-a,\ X(x_i)=W_i$とすれば
$(E\wedge X)^\sharp dh=hX-X(h)E=V$である。
よって$\widetilde P=\Pi-E\wedge X$は$h$をCasimirとする。

別のこのようなliftとの差が$E\wedge Y$なら$hY-Y(h)E=0$であり、
多項式性から$Y$はEuler場の定数倍、従って差は零である。
射影Jacobi条件から$E\wedge[\widetilde P,\widetilde P]=0$。
$\widetilde P$は$\partial_h$成分を持たないため、この等式の$\partial_h$成分は
$h[\widetilde P,\widetilde P]=0$を与え、真のJacobi恒等式が従う。
$h=1$で$P$を得る。

逆に任意の二次以下の$P$に対し
$\widetilde P_{ij}=h^2P_{ij}(x/h),\ \widetilde P_{hi}=0$とすれば斉次二次Poisson liftを得る。
Jacobiは$h\ne0$で確認でき、多項式恒等式として全域に延長する。
$\wedge^4T_{\PP^4}=\cO(5)$なので、そのPfaffian切断は式\eqref{eq:marked-divisor}である。
最高次項が非零であることから$H$の重複度も従う。
\end{proof}

\begin{proposition}[標識を忘れる際の有限性]\label{prop:finite-mark}
一般階数四の超平面因子を持つ構造の空間に対し、対$([\pi],H)$の標識空間からの忘却写像は有限である。
各幾何学的ファイバーは高々五点である。
標識された構造の射影同型類は、アフィン二次Poisson構造のアフィン座標変換と非零スカラー倍による同型類に対応する。
\end{proposition}
\begin{proof}
$h\mid\Pf(\pi)$は$\pi$と$[h]$に関する閉incidence条件であり、超平面のパラメータ空間は射影的である。
したがって忘却写像は固有である。非零五次式の一次既約因子は高々五つなので準有限でもあり、有限となる。
$H$の$\PGL_5$安定化群は$\mathbb A^4$上のアフィン変換群であり、定理\ref{thm:marked-affine}の一意性から最後の主張が従う。
\end{proof}
標識の選び直しで同じ成分を重複計数し得ることは残る。
この有限性はrank二locusには主張しない。

\subsection{全超平面型に共通のmodular--Pfaffian復元}
$P=\sum_{i<j}p_{ij}\partial_i\wedge\partial_j$、
$Z^j=\sum_i\partial_ip_{ij}$、$\mu=dx_1\wedge\cdots\wedge dx_4$とする。
$P^\sharp(df)^j=\sum_i p_{ij}\partial_if$という規約を用いる。

\begin{proposition}[modular原始関数による全源の判定]\label{prop:modular-primitive}
任意の多項式二ベクトル$P$に対し、$\alpha=\iota_P\mu$、$g=\Pf(P)$と置く。
$J^{ijk}=\sum_l(p_{il}\partial_lp_{jk}+p_{jl}\partial_lp_{ki}+p_{kl}\partial_lp_{ij})$
および$J=\sum_{i<j<k}J^{ijk}\partial_i\wedge\partial_j\wedge\partial_k$とすると
\begin{equation}
 d\alpha=-\iota_{\diver P}\mu,\qquad
 dg+\iota_{\diver P}\alpha=\iota_J\mu.\label{eq:primitive-identity}
\end{equation}
従って、固定した$Z$と$f$（ここでは$f=0$も許す）に対し、全源は正確に
\begin{equation}
 d\alpha=-\iota_Z\mu,\qquad
 \iota_Z\alpha=-df,\qquad
 \alpha\wedge\alpha/2=f\mu                         \label{eq:primitive-source}
\end{equation}
である。初めの二条件は$\alpha$に関してアフィン線形で、最後だけが二次である。
解では$\diver Z=0$、$Z(f)=0$、$\mathcal L_Z\alpha=0$、$[Z,P]=0$が従う。
\end{proposition}
\begin{proof}
$\alpha=p_{12}dx_3\wedge dx_4-p_{13}dx_2\wedge dx_4
+p_{14}dx_2\wedge dx_3+p_{23}dx_1\wedge dx_4
-p_{24}dx_1\wedge dx_3+p_{34}dx_1\wedge dx_2$を微分すると第一式を得る。
第二式の$dx_4$係数は
$\partial_4(p_{12}p_{34}-p_{13}p_{24}+p_{14}p_{23})
+Z^1p_{23}-Z^2p_{13}+Z^3p_{12}=J^{123}$である。
添字の交代置換で他の三係数も得る。従ってPfaffianが零でもこの恒等式は成立し、
$J=0$との同値から全源表示を得る。
第一条件を微分して$\diver Z=0$、第二条件を$Z$で縮約して$Z(f)=0$を得る。
Cartan公式から
$\mathcal L_Z\alpha=d\iota_Z\alpha+\iota_Zd\alpha=0$であり、
$\mathcal L_Z\mu=0$と外積双対の一意性から$[Z,P]=0$となる。
\end{proof}

この表示は単に固定因子を線形に書き直すだけでなく、以下で単純Hと重複Hの
\emph{入口の証明}を同時に与える。$f=0$で必要なのは$\iota_Z\alpha=0$である。
弱い式$P^\sharp df=fZ$だけでは代用できない。

\begin{theorem}[固定した$f,Z$からの全点復元]\label{thm:modular-reconstruct}
$f\ne0$を次数四以下の多項式、$Z$をアフィン線形場として固定する。
$f=\Pf(P)$、$Z=\diver P$を満たす二次以下のPoisson構造の集合は、
アフィン線形空間
\begin{equation}
 \mathcal A_{f,Z}=\{P:\deg P\le2,\ \diver P=Z,\ P^\sharp df=fZ\}
                                                               \label{eq:marked-linear}
\end{equation}
に二次条件$\Pf(P)=f$を課した集合と正確に一致する。
従って固定した$f,Z$に対して、追加のJacobi方程式は不要である。
空でない場合には$\diver Z=0$、$Z(f)=0$、$[Z,P]=0$が成り立つ。
\end{theorem}
\begin{proof}
$\alpha=\iota_P\mu$とする。成分表示から
$d\alpha=-\iota_{\diver P}\mu$、
$d f\wedge\alpha=-\iota_{P^\sharp df}\mu$である。
$\Pf(P)=f\ne0$のとき、逆形式の閉性は
\[
 f\,d\alpha-df\wedge\alpha=0
 \quad\Longleftrightarrow\quad P^\sharp df=f\diver P
\]
である。非退化双ベクトルのJacobi恒等式は逆形式の閉性と同値であるから、結論を得る。
式\eqref{eq:marked-linear}が$P$についてアフィン線形、Pfaffianが二次であることも明らかである。
$\diver Z=0$は交代係数への二重微分の相殺から従う。
$df$で式\eqref{eq:marked-linear}を評価すると$fZ(f)=0$、従って$Z(f)=0$。
また$f\ne0$上で$Z=P^\sharp(d\log f)$は局所Hamiltonian場なので$[Z,P]=0$。
これは正則テンソルの等式として全域に成り立つ。
\end{proof}

\begin{corollary}\label{cor:unimodular}
一般階数四の標識超平面模型において
\[
 Z=0\quad\Longleftrightarrow\quad f\in\CC^\times
 \quad\Longleftrightarrow\quad D_\pi=5H.
\]
\end{corollary}
\begin{proof}
$Z=0$なら$P^\sharp df=0$。非退化開上で$df=0$なので$f$は定数である。
逆は同じ恒等式と$f\ne0$から従い、最後の同値は式\eqref{eq:marked-divisor}である。
\end{proof}

この定理は、すべての$f,Z$を既に分類したという主張ではない。
$f,Z$を同時に未知数に戻すと式\eqref{eq:marked-linear}はもはや$P$だけの線形問題ではない。
スカラー倍では$(P,f,Z)\mapsto(cP,c^2f,cZ)$を同時に保つ。
$Z$の半単純・冪零分解と$Z(f)=0$を先に利用し、
重複超平面、$H+Q_4$、複数超平面の場合に共通の復元定理として適用することが狙いである。
とくに$\deg f\le3$の枝でのみ$\Pf(P_2)=0$が従い、これを単純超平面へ誤用してはならない。

\begin{proposition}[共通の有限共鳴層別]\label{prop:resonance}
$Z$の半単純・冪零分解を$Z=S+N$とする。
アフィン座標を選んで$S=\sum_{i=1}^4\lambda_i x_i\partial_i$、$[S,N]=0$とできる。
Poisson解$P$は$[S,P]=[N,P]=0$を満たす。
$x_0=1,\lambda_0=0$と書けば、$P_{ij}$に現れる単項式$x_kx_l$には
\begin{equation}
 \lambda_k+\lambda_l=\lambda_i+\lambda_j
 \quad(1\le i<j\le4,\ 0\le k\le l\le4)              \label{eq:marked-resonance}
\end{equation}
が必要である。固有値一致条件も加えると、$\sum\lambda_i=0$の空間は
許容単項式と冪零blockの組合せが一定な有限個の構成可能な層に分かれる。
固有値比そのものが有限個になるとは主張しない。
\end{proposition}
\begin{proof}
アフィン線形場を五次元拡大空間の行列として表示し、$h$を保つJordan分解を取る。
半単純部分の固定点へ平行移動し、その固有空間で線形座標を選べば上の形になる。
次数二以下の多項式双ベクトルの有限次元表現において$\operatorname{ad}S$は半単純、
$\operatorname{ad}N$は冪零で可換である。$\operatorname{ad}S$の固有値$\nu\ne0$の空間では
$\nu I+\operatorname{ad}N$が可逆なので、$[Z,P]=0$からその成分は零となる。
残る零固有空間では$[N,P]=0$となる。式\eqref{eq:marked-resonance}は単項式の重み比較である。
有限個の線形等式の成立・不成立による層別は有限であり、
各重複固有空間のJordan分割も有限個である。
translationの余核とその零・非零の分岐はこの後にも保持し、無断で捨てない。
\end{proof}

\begin{remark}[rank二での失敗を明示する]
$P=\partial_1\wedge(\partial_2+x_1\partial_3)$は$f=0$を満たすので
$P^\sharp df=f\diver P$は恒等的に成立するが、
$\{x_1,\{x_2,x_3\}\}+\{x_2,\{x_3,x_1\}\}+\{x_3,\{x_1,x_2\}\}=1$でありPoissonではない。
$f\ne0$は定理の本質的仮定である。
\end{remark}

\subsection{二つの旧入口を置き換え、第三の配置で確認する}
$P=C+L+Q$に対応して$\alpha=\alpha_0+\alpha_1+\alpha_2$、
$f=\sum_{k=0}^4 f_k$と係数次数で分解する。
命題\ref{prop:modular-primitive}は、全低次項を残したまま
\begin{align}
 d\alpha_1&=-\iota_c\mu,&d\alpha_2&=-\iota_{Ax}\mu,\label{eq:degree-divergence}\\
 df_k&=-\iota_c\alpha_{k-1}-\iota_{Ax}\alpha_{k-2}
       &&(1\le k\le4)                              \label{eq:degree-primitive}
\end{align}
と$\alpha^2/2=f\mu$を与える。範囲外の$\alpha_j$は零とする。
特に最高式は
\begin{equation}
 d\Pf(Q)=-\iota_{Ax}\alpha_2.\label{eq:highest-primitive}
\end{equation}
これは$Q$自体に対するJacobi式であり、$C,L$を解から消してよいという意味ではない。

\paragraph{単純Hと既約四次.}
$D=hF_4$では$\deg f=4$であり、$K=\Pf(Q)\ne0$。
式\eqref{eq:highest-primitive}により$A=0$は不可能である。
実際$A=0$なら$dK=0$となり、非零四次斉次式に反する。
一方$\iota_{Ax}\alpha_2$は零ではないので、$Ax\wedge Q=0$という
重複H専用の制約を課してはならない。
全共鳴は命題\ref{prop:resonance}から得る。
ここでなお必要なのは、無限遠$K$の閉有理一形式による還元、三つの一般wall上の全解排除、
三十七wall・十九低次wallの有限列挙と、その後の全源帰属である。
共通共鳴だけから二十四代表や各成分の極大性は従わない。

\paragraph{重複H.}
$h^2\mid D$は$\deg f\le3$、すなわち$K=0$と同値である。
式\eqref{eq:highest-primitive}から$\iota_{Ax}\alpha_2=0$、従って
$(Ax)\wedge Q=0$を得る。
これが従来別々に導いた最高部のKoszul入口である。
さらに$f_3,f_2,f_1$の式を全て残すため、最高部だけの支持への置換にはならない。
$\rank A\ge3$でのKoszul完全性、各Jordan型の残二次条件と全零補集合、
実変形による帰属は従来通り必要である。

\paragraph{共通の正逆とtranslation.}
両入口の逆は同じで、式\eqref{eq:primitive-source}から$P$を外積双対で戻し、
$h^2P(x/h)$を斉次化するだけである。座標多項式や源パラメータを反転しない。
アフィン変更$y=Mx+b$では
$P'=M P(M^{-1}(y-b))M^{\mathsf T}$、$Z'=MZ(M^{-1}(y-b))$、
$f'=\det(M)f(M^{-1}(y-b))$で、$\det(M)\ne0$のみを用いる。
特に平行移動後の定数modular項は$c-Ab$（$M=I$）であり、
$c$を消せるのはその$\operatorname{im}A$成分だけである。
残る$\operatorname{coker}A$、異なるblock長のtop、共鳴wall、$C$の全階数を保持する。

\paragraph{第三の配置：二つの単純Hと三次因子.}
$D=hxF_3$、$F_3$既約、$h,x$相異なる場合、$h$を標識した模型では
$P^{xi}$は$x$で割れ、全源は
\[
 P=x\partial_x\wedge V+W,\qquad
 V=(By+cx+d)\cdot\partial_y,\quad \deg W\le2.
\]
逆は$P^{xi}/x$の\emph{係数の取り出し}とvertical blockの取り出しであり、
開集合$x\ne0$への局所化ではない。全45係数の$W$と$c,d$を保持する。
$F_3|_{h=1}=F$とすると$f=xF$であり、これは通常次数四を持つ単純Hの入口である。
$Z=-x\operatorname{tr}(B)\partial_x+V+x\partial_xV+\diver_y W$を
式\eqref{eq:primitive-source}に代入したJacobi条件は正確に
\[
 [V,W]+xV\wedge\partial_xV=0,\qquad
 [W,W]/2=xV\wedge\partial_xW
\]
となる。これは元の四Jacobiを$x$を含む三組と残る一組へ分けたものであり、
共通命題の正逆が第三配置でも全点を保つ。
$F=\iota_V\iota_W(dy_1\wedge dy_2\wedge dy_3)$の既約性と次数三、
ならびに残る全sourceの帰属は別に課す。

\paragraph{置換の正確な範囲.}
主分類稿の「重複超平面とモジュラー場」の導入と
補題modular-koszulの$Ax\wedge Q=0$の導出、
「単純超平面と既約四次因子の有限入口」の標識・発散・Jordan共鳴の導入を、
本節の同じ証明へ置き換える。
原証明は\path{HYPERPLANE_QUARTIC_COFACTOR_ENTRANCE.md}第1・3節、
\path{KNOWN142_SEVEN_CANDIDATE_WHOLE_POISSON_COVER_AUDIT.md}第4節とそこからの全源である。
第三配置の独自の三次式認識は\path{TWO_HYPERPLANES_CUBIC_SUPPORT.md}第1節に残る。
閉形式の留数による認識、有限wall列挙、個別源の成分性と非包含を削除していない。
本節の通常の証明の新規性は未確認であり、既存147成分全体の独立検証を主張しない。

\subsection{定数modular場の全源を、積分定数まで復元する}
重複Hの中でも$Z=-\partial_t$の全源は、三座標$u=(u_1,u_2,u_3)$を用いて
\begin{equation}
 P=\partial_t\wedge Y+\operatorname{Jac}_u(f),\qquad
 \diver_uY=1,\quad Y(f)=f,\quad \deg Y\le2,\quad\deg f\le3
                                                        \label{eq:constant-whole}
\end{equation}
と必要十分に書ける。ここで$\operatorname{Jac}(f)^{23}=f_{u_1}$、循環的に定める。
\begin{proof}
$[Z,P]=0$から全係数は$t$に依存しない。
縦二ベクトルを$W$、$\Omega=du_1\wedge du_2\wedge du_3$、
$\beta=\iota_W\Omega$とすれば
$\alpha=\iota_Y\Omega+dt\wedge\beta$である。
式\eqref{eq:primitive-source}の第二式は$\beta=df$を与え、
第一式は$\diver Y=1$、最後は$\beta(Y)=f$、すなわち$Y(f)=f$を与える。
逆にこれらを満たせば三条件が全て成立する。
従って$f$の加法定数も自由に捨てられず、$Y(f)=f$によって固定される。
この正逆は$f_3$の被約性、磁気係数の非零性を仮定せず、全lower sourceを保つ。
rank四を扱うときは$f\ne0$だけを別に課す。
\end{proof}


% Shared structural proofs; no invariant hyperplane or affine quadratic chart.
\section{二次＋三次のcofactor復元と平面型の磁気商}\label{sec:quadric-reconstruction}
この節では超平面因子を仮定せず、五変数の斉次liftを用いる。
前半は二次因子の階数三・四・五を統一し、後半は平面型を因子分割によらず復元する。
後半で得る磁気商の次元は、単に固定因子の線形表示を与えるだけでなく、
従来個別に計算した磁気shearの階数と零補集合、軌道次元の違いを説明する。

\subsection{発散零liftとPfaffianの原始関数}
五変数$z=(z_0,\ldots,z_4)$のEuler場を$E$、定数体積を$\Omega$とする。
交代二次行列$P$の発散を$Z^j=\sum_i\partial_iP^{ij}$とし、
\[
 \mathcal C(P)=P-\tfrac15 E\wedge Z
\]
と置く。$\mathcal C$は同じ射影双ベクトルを表す発散零liftへの忠実な射影である。
実際$\operatorname{div}Z=0$と
$\operatorname{div}(E\wedge Y)=5Y-(\operatorname{tr}Y)E$から発散零となる。
発散零な二つのliftの差が$E\wedge Y$なら$5Y=(\operatorname{tr}Y)E$なので差は零であり、一意である。
この操作とその射影空間への降下はGL五変数の実座標変更と可換である。

\begin{proposition}[斉次Pfaffianの全Jacobi判定]\label{prop:cone-pfaffian-primitive}
$\Pi$を発散零の斉次二次二ベクトルとし、
\[
 \alpha=\iota_\Pi\Omega,\qquad
 \beta=\iota_{\Pi\wedge\Pi/2}\Omega,\qquad
 D=\iota_E\beta
\]
と定める。$\alpha$は閉三形式、$\beta$は係数次数四の一形式、$D$は五次式である。
このとき次は同値である。
\[
 \text{射影Jacobi}\quad\Longleftrightarrow\quad
 [\Pi,\Pi]=0\quad\Longleftrightarrow\quad
 d\beta=0\quad\Longleftrightarrow\quad 5\beta=dD.
\]
$D\ne0$のとき$V(D)$はactual退化因子で、$\Pi^\sharp dD=0$である。
$D$をJacobi前に計算できても、条件を課す前は単なるpotentialである。
\end{proposition}
\begin{proof}
$J=[\Pi,\Pi]/2$とする。成分の反対称化で得る恒等式は
\[
 d\iota_{\Pi^2/2}\Omega+
 \iota_{\operatorname{div}\Pi}\iota_\Pi\Omega=\iota_J\Omega.
\]
従って発散零なら$d\beta=\iota_J\Omega$。
また$\mathcal L_E\beta=5\beta$から
$dD=5\beta-\iota_Ed\beta$である。
射影Jacobiは$E\wedge J=0$、すなわち$\iota_Ed\beta=0$なので$5\beta=dD$となり、
もう一度微分して$d\beta=0$、従って$J=0$を得る。逆は明らかである。
奇数次交代行列の符号付きPfaffian列はその行列の核に入るから、$\Pi^\sharp\beta=0$。
$D$は射影外積$\pi^2/2$の切断を表し、最後の主張も従う。
\end{proof}

\begin{theorem}[全二次＋三次の標識付き復元]\label{thm:quadric-cofactor-reconstruct}
actual因子が$D=qc$で、$q,c$はそれぞれ既約二次・既約三次とする。
唯一の発散零liftに対し、唯一の非零線形場$X$が存在して
\begin{equation}
 \Pi^\sharp dq=qX,\quad Xq=Xc=0,\quad
 \operatorname{tr}X=0,\quad [X,\Pi]=0.\label{eq:quadric-cofactor}
\end{equation}
射影上の因子由来の対称性
$\pi^\sharp(\frac12d\log q-\frac13d\log c)$は$\frac56X$の降下である。
逆に$q,c,X$を固定した全源は、斉次二次$\Pi$について正確に
\begin{equation}
 \operatorname{div}\Pi=0,\qquad \Pi^\sharp dq=qX,\qquad
 \iota_{\Pi^2/2}\Omega=\tfrac15d(qc)                 \label{eq:quadric-whole-source}
\end{equation}
で復元される。最初の二条件はアフィン線形、最後は二次であり、追加Jacobiは不要である。
不変性$[X,\Pi]=0$を線形源の構成時に課しても全解は変わらない。
\end{theorem}
\begin{proof}
$\Pi^\sharp d(qc)=0$と$\gcd(q,c)=1$から$q$は$\Pi^\sharp dq$の全係数を割る。
商は次数一の場$X$で、一意である。同じ式から$\Pi^\sharp dc=-cX$。
反対称性から$Xq=0$、$\{q,qc\}=0$から$Xc=0$を得る。
$X=0$なら$q,c$がともにCasimirとなるが、次数二・三の相異なる既約式は代数的に独立である。
従って一般五次元で核一次元のrank四二ベクトルには不可能である。
Hamiltonian場$qX$の発散は発散零liftでは零なので
$q\operatorname{tr}X+X(q)=0$からtrace零となる。
$q\ne0$上で$X$は局所Hamiltonian場$\Pi^\sharp d\log q$であり、$[X,\Pi]=0$は全域へ延長する。
対数一形式の係数$1/2+1/3=5/6$が対称性の比例係数を与える。
逆は命題\ref{prop:cone-pfaffian-primitive}による。
最後の二次条件は$D=qc\ne0$を正確に規格化するのでrank四も保証される。
全条件が多項式であり、$q$の非零開にだけある解を多項式全源と取り違えていない。
\end{proof}

補集合上の閉逆二形式を$\omega$、$\eta=\frac12d\log q-\frac13d\log c$とすると、
本稿のsharp規約では$\iota_{(5/6)\overline X}\omega=-\eta$である。
従ってこの復元は、因子由来の閉一形式に対応する対称性を持つ閉逆形式の復元でもある。
固定$D$での閉有理二形式の線形・Pfaffian二次表示と式\eqref{eq:quadric-whole-source}は、
一般rank四の同じ全構造を表す。対称性の存在だけで閉形式の全fiberを解いたとはしない。

因子を非零定数倍しても$X$は変わらない。tensorを$\lambda$倍すると
$(\Pi,D,X)$は$(\lambda\Pi,\lambda^2D,\lambda X)$となる。
本節の$1/5$は$D=\iota_E\beta$というactual切断の規約による。
旧sourceのpotential$\iota_E\beta/5$を使う規約とは非零定数五だけ異なり、因子・支持は同じである。

\subsection{radicalの余核と全31代表を残す}
$q$の非退化空間とradicalを$U,R$とすると、式\eqref{eq:quadric-cofactor}から
\[
 X=\begin{pmatrix}A&0\\ B&C\end{pmatrix},\qquad
 A\in\mathfrak{so}(q_U),\quad C\in\mathfrak{sl}(R).
\]
特にrank四では$C=0$であり、radical上の一般線形場を許すのは誤りである。
実shear$r'=r+Tu$とその逆は全域で定義され、
\begin{equation}
 B\longmapsto B+TA-CT.\label{eq:quadric-sylvester}
\end{equation}
従って残るのはSylvester写像$T\mapsto TA-CT$の余核である。
スペクトルが交わらない層ではこの写像は可逆だが、共鳴層では余核を保持し、そのcentralizer作用で分ける。
例えばrank三で$A=a\partial_b-b\partial_c,C=r\partial_s$の場合、全shearで
$Xr=pc,Xs=r+qc$にでき、残る中心化群は
$(p,q)\mapsto(\lambda p,\mu p+\lambda q)$、$\lambda\ne0$と作用する。
従って$p\ne0$、$p=0,q\ne0$、$p=q=0$の三つを全て残す。
これは共鳴の余核を消してよいという一般論ではない。

半単純スペクトルは$(0,a,-a,b,-b)$である。
$ab\ne0$で絶対比が二以外なら、weight零の全三次単項式は零固有座標で割れる。
従って既約$c$は正則の絶対比二、非正則の零固有空間、純冪零の各型へ帰着する。
rank三・四の全shear余核とcentralizerの分類は、元の全表
\path{DEGENERATE_QUADRIC_CUBIC_COFACTOR_REDUCTION.md}第3--5節を保持する。
\begin{center}\small
\begin{tabular}{c|l|r|l}
rank $q$&$X$の群&代表数&なお使う元の全支持入力\\\hline
4&正則半単純&1&CUBIC\_REGULAR\_SUPPORT\\
4&非正則半単純&3&CUBIC\_SEMISIMPLE\_SUPPORT\\
4&純冪零$A\ne0$&5&Q4\_NIL\_SUPPORT\\
4&$A=0$&2&CUBIC\_REGULAR\_SUPPORT第3節\\
3&正則半単純&2&CUBIC\_REGULAR\_SUPPORT\\
3&非正則半単純&6&CUBIC\_SEMISIMPLE\_SUPPORT\\
3&純冪零$A\ne0$&5&Q3\_ANIL3\_SUPPORT\\
3&$A=0$&7&Q3\_AZERO\_SUPPORT
\end{tabular}
\end{center}
表の接頭辞は\path{DEGENERATE_QUADRIC_}、末尾は\path{.md}である。
rank五の全Jordan型は\path{SMOOTH_QUADRIC_CUBIC_SUPPORT.md}第1・6・7節にある。
式\eqref{eq:quadric-whole-source}は全型に同じ正逆を与えるが、各源の既約$c$条件・零補集合・実変形・
成分帰属はこれらの元証明を必要とする。31代表を31成分と数えない。

\subsection{被約平面五次の完全復元は因子分割によらない}
以下では$S=\CC[a_0,a_1,a_2]$、$F\in S_5$を被約とし、頂点方向を$r,s$とする。
$E_0=\sum a_i\partial_{a_i}$、$\Omega_0=da_0\wedge da_1\wedge da_2$と置く。
\begin{theorem}[全平面五次のSaito源]\label{thm:reduced-plane-whole-source}
actual因子が$F$であるrank四構造が存在すれば、$F$はexponents$(2,2)$のfree平面五次である。
その全発散零源は、二次gradient syzygyの全frame$U,V$と任意の$R\in S_2$を用いた
\begin{equation}
 \Pi_R=U\wedge\partial_r+V\wedge\partial_s+
 \{-(\operatorname{div}V)r+(\operatorname{div}U)s+R\}
       \partial_r\wedge\partial_s                         \label{eq:plane-whole-source}
\end{equation}
で尽くされる。逆にこの自由基底を持つ全被約$F$と全$R$がrank四Poisson源を与える。
$\det(E_0,U,V)=\kappa F$、$\kappa\ne0$ならactual切断は$-\kappa F$で、$R$に依存しない。
定数・一次syzygyは無く、二次syzygyは正確に二次元、全Tjurina数は十二である。
\end{theorem}
\begin{proof}
命題\ref{prop:cone-pfaffian-primitive}から$\Pi^\sharp dF=0$。
baseの交代三行列をvector$b$で表せば$b\times\nabla F=0$である。
被約性とEuler式から$\nabla F$の係数の最大公約式は一。
従ってUFDの各primeのvaluationを比べれば$b$は$\nabla F$の多項式倍だが、
次数二と四の差から$b=0$。頂点依存を仮定せずにbase blockが零となる。
全二次mixed列を$U,V$、残る項を$B$とすれば
$U(F)=V(F)=0$、$D=-\det(E_0,U,V)$である。
$r=s=0$の$U_0,V_0$と$E_0$の行列式は非零定数倍の$F$。
任意のlogarithmic場のCramer分子は、各reduced factorの一般滑らかな点で零となる。
従って$F$で割れ、$\operatorname{Der}(-\log F)=SE_0\oplus SU_0\oplus SV_0$。
$E_0(F)=5F$なので$AR(F)=SU_0\oplus SV_0$も得る。
よって低次syzygyは無く、$U,V$の正頂点次数の係数は全て零である。

行列式のLie微分から$[U,V]=(\operatorname{div}V)U-(\operatorname{div}U)V$。
Jacobiのbase,$r,s$成分は$[U,V]+B_rU+B_sV=0$なので、
$B_r=-\operatorname{div}V,B_s=\operatorname{div}U$となり、式\eqref{eq:plane-whole-source}を得る。
逆にこの式では全Jacobiと発散零が成立し、Pfaffianは上の$-\kappa F$に等しい。
全frameの正逆は定数GL二変数だけであり、六係数の$R$を捨てない。
自由syzygy分解からJacobian idealの解像は
$0\to S(-6)^2\to S(-4)^3\to J_F\to0$。
$S/J_F$のHilbert分子$1-3t^4+2t^6$の二階微分の半分は十二であり、全Tjurina数を与える。
\end{proof}

\subsection{磁気商は因子の留数によって測られる}
$F=\prod_{i=1}^{\nu}f_i$、$d_i=\deg f_i$とし、因子は相異なる既約式とする。
全域で可逆なshear$r'=r+\ell(a),s'=s+m(a)$、$\ell,m\in S_1$は、正確に
\[
 R\longmapsto R+\delta_F(\ell,m),\qquad
 \delta_F(\ell,m)=\operatorname{div}(\ell V-mU)
\]
と作用する。逆は$\ell,m$を引くだけであり、一次近似の作用ではない。
\begin{theorem}[磁気shearの核と因子留数]\label{thm:magnetic-residue-quotient}
式\eqref{eq:plane-whole-source}の全被約$F$について、自然な同型
\begin{equation}
 \ker\delta_F\ \simeq\
 \{(\lambda_i)\in\CC^\nu:\textstyle\sum_i d_i\lambda_i=0\}
                                                        \label{eq:magnetic-residues}
\end{equation}
がある。従って
\[
 \dim\ker\delta_F=\dim\operatorname{coker}\delta_F=\nu-1,
 \qquad\rank\delta_F=7-\nu.
\]
核のshearは因子由来の可換Poisson対称性そのものである。
磁気商$\mathcal M_F=S_2/\operatorname{im}\delta_F$の零点を捨ててはならない。
\end{theorem}
\begin{proof}
核の$(\ell,m)$に$W=\ell V-mU$、$\theta=\iota_W\Omega_0$を対応させる。
$\operatorname{div}W=0,W(F)=0$だから$\theta/F$は閉二形式、Euler重み零である。
よって$\eta=\iota_{E_0}(\theta/F)$は閉・水平・重み零の一形式で、$F$に沿って高々単純極を持つ。
各因子の一般滑らかな点で閉性は定数留数$\lambda_i$を与える。
一般直線に制限した留数定理から$\sum_i d_i\lambda_i=0$。
$\sum_i\lambda_i df_i/f_i$を引くと全divisorial polesが消え、
余次元二延長と$H^0(\PP^2,\Omega^1)=0$から差は零となる。
$\eta=0$なら$E_0,W$は関数体上平行だが$E_0(F)=5F,W(F)=0$だから$W=0$。
自由基底の独立性から$\ell=m=0$であり、単射を得る。

逆にこの次数関係を満たす対数一形式$\eta=\sum_i\lambda_i df_i/f_i$を取る。
$\gamma=F\eta$は多項式一形式で、水平かつEuler重み五。
\[
 \iota_W\Omega_0=\tfrac15d\gamma
\]
で一意な三次場$W$を定めると、$\operatorname{div}W=0$、
$dF\wedge d\gamma=dF\wedge dF\wedge\eta=0$から$W(F)=0$。
Cartan公式により$\iota_{E_0}\iota_W\Omega_0=\gamma$である。
自由分解$AR(F)_3=S_1U\oplus S_1V$から一意な$\ell,m$を回収でき、全射となる。
定義域$S_1^2$と値域$S_2$はともに六次元なので、余核の次元も同じである。

最後に因子の対数場は
$\Pi_R^\sharp\eta=\eta(U)\partial_r+\eta(V)\partial_s$。
各$U,V$が全因子の対数場なので係数は一次式である。
$W_\eta=\eta(U)V-\eta(V)U$について
$\iota_{E_0}\iota_{W_\eta}\Omega_0=\kappa F\eta$となるから、
非零定数$\kappa$を除き上の同型と一致する。
\end{proof}

\begin{proposition}[磁気零境界と軌道次元]\label{prop:plane-magnetic-orbits}
固定$F$の全frameをGL二変数で合わせると、残るshear商は$\mathcal M_F$であり、
共通fiber scalarとtensor scaleの組合せは$[R]\mapsto t[R]$と作用する。
base曲線の自己同型群は有限である。
従って$[R]=0$の軌道次元は$24-\nu$、$[R]\ne0$の軌道次元は$25-\nu$である。
特に$\nu=1$では磁気項は全消去でき、一つの23次元軌道となる。
$\nu=2$では非零磁気商が一つの23次元軌道、零商が22次元の境界軌道となる。
\end{proposition}
\begin{proof}
$\operatorname{Der}(-\log F)_1=\CC E_0$なのでPGL三変数の曲線安定化群は零次元、従って有限。
本質三変数空間は$F$から内在的に決まる。
tensor安定化群の連結部分はbaseを固定し、fiberに$M(r,s)^t+Ca$と作用する。
anchorの独立性から射影tensorを保存する$M$はscalarであり、shearの自由度は
式\eqref{eq:magnetic-residues}の$\nu-1$である。
零商ではさらにfiber scalarの一次元が残り、非零商ではその連結自由度は無い。
$\dim\PGL_5=24$から結論を得る。$\nu=2$の零商は$R_t=tR$の実極限で、全tでactual因子は同じである。
軌道次元やこの極限だけから全Poisson空間の成分性を主張していない。
\end{proof}

\paragraph{実際に置換する旧証明と保持する計算.}
滑らかな二次因子と特異二次因子のcofactor・trace・不変性の二つの入口を
定理\ref{thm:quadric-cofactor-reconstruct}へ統合する。
平面型の完全源は二H＋三次、rank三Q＋三次、既約五次の三箇所で
定理\ref{thm:reduced-plane-whole-source}へ統合する。
Q＋三次の六源の磁気shear階数五と、既約五次の階数六は
定理\ref{thm:magnetic-residue-quotient}の$\nu=2,1$の同じ証明になる。
二H＋三次では$\nu=3$なので商次元二であり、因子数を無視して全磁気消去を主張しない。
一方、全31場の有限分類、各源の既約性・零補集合、六つの平面配置の尽くし、
周囲の二次・三次障害と成分相違、非平面型の実変形・S61帰属は旧証明を保持する。
因子衝突で被約性が失われる境界へ、この留数次元式を延長しない。
固定した因子数のfree被約層では磁気写像は定数階数で、minor chart上に核・像・余核の束を作れるが、
一般層への制限で特殊底成分を捨てる意味ではない。
以上の統合の文献上の新規性は未確認であり、全147成分の独立検証ではない。


\subsection{既約五次に残る導手の認識}
第三の入口のうち非正規の場合は、正規化の導手支持と、その上のrestriction rank零・二を使う。
現行のケース別証明では、一般平面切断の算術種数六から得る
$\sum_i\delta_i\deg S_i=6-g$という有界性と、
導手上の留数・余法線の一次記号が重要な入力になっている。
これらの全支持認識を本稿の基本定理で証明済みとはしない。
正規の場合の楕円認識、非正規の場合の導手認識を、それぞれ独立した構造定理として再構成する。

\section{固定因子の復元を相対化し、変形の持上げを分ける}
$X=\PP^4$、$L=\cO(5)$とし、
\[
 \mathbb E=H^0(X,\Omega_X^2\otimes L),\quad
 \mathbb B=H^0(X,L),\quad
 \mathbb C_3=H^0(X,\Omega_X^3\otimes L^2)
\]
と置く。$\mathbb C_3$は係数空間の記号であり、体$\CC$とは区別する。
$d_F:\mathbb E\to\mathbb C_3$は前節で定義した線形写像である。

\begin{proposition}[相対復元]\label{prop:relative}
$\PP(\mathbb E)\times\PP(\mathbb B)$の局所閉集合
\[
 \mathcal I^\circ=
 \{([\alpha],[F]):d_F\alpha=0,\ q(\alpha)\in\CC^\times F\}
\]
は、一般階数四のPoisson空間から退化因子への写像
$\delta:[\pi]\mapsto[\pi^2/2]$のグラフと同一視できる。
$[F]$のパラメータ層で$\rank d_F$が一定なら、$\ker d_F$はベクトル束をなし、
その射影束内でPfaffian二次条件を課して$\mathcal I^\circ$を得る。
\end{proposition}
\begin{proof}
最初の主張は定理\ref{thm:fibre}の全$F$での適用であり、$q(\alpha)\ne0$なら$[F]$は一意である。
$F$のスカラー倍は$d_F$も同じ倍にするため、そのkernelは$[F]$だけで決まる。
定数階数の層では可逆minorを使って線形系を解け、kernelは局所自由になる。
Pfaffianの$\mathbb B/\CC F$への像の消失は斉次二次条件である。
主張は被約台の表示であり、全非被約基底上の変形関手の同値まで含めない。
\end{proof}

\begin{proposition}[一次・二次の持上げ方程式]\label{prop:lifting}
$q(\alpha)=F\ne0$、$d_F\alpha=0$と規格化する。
$\alpha_t=\alpha+ta+t^2b$、$F_t=F+t\varphi+t^2\psi$について、
Poisson構造とその退化因子の変形条件は、$t^3$を法として次の二段階で与えられる。
\begin{align}
 \alpha\wedge a&=\varphi,&
 d_Fa&=d\varphi\wedge\alpha-\varphi\,d\alpha; \tag{5.1}\label{eq:firstlift}\\
 \alpha\wedge b&=\psi-\tfrac12a\wedge a,&
 d_Fb&=d\psi\wedge\alpha-\psi\,d\alpha+d\varphi\wedge a-\varphi\,da.\tag{5.2}\label{eq:secondlift}
\end{align}
特に固定$F$の一次変形空間は
$\{a\in\ker d_F:\alpha\wedge a=0\}$である。
\end{proposition}
\begin{proof}
$q(\alpha_t)=F_t$と$F_t d\alpha_t-dF_t\wedge\alpha_t=0$の各次数を展開する。
$F\ne0$の開集合上で逆形式を使えばJacobi条件と同値である。
元の$F$が非零の開集合への制限は、$t^3=0$の係数環上でも正則テンソルについて単射なので、
そこでのJacobi等式は全域の等式を与える。
\end{proof}

これにより$\Delta_\alpha(a)=(\alpha\wedge a,d_Fa)$と置けば、
指定された因子変形$\varphi$の持上げの障害は
\[
 \bigl[(\varphi,d\varphi\wedge\alpha-\varphi d\alpha)\bigr]
       \in\operatorname{coker}\Delta_\alpha
\]
である。ここで示したのは指定された方程式の線形な可解性であって、
全高次変形のunobstructednessではない。
因子を変形できること、Poisson構造まで持ち上がること、持上げた族が他の成分へ入ることを分離できる。
この区別を保った上で、次節では実際の二つの族で補正の必要性を比較する。

% Shared applications; all coefficients below use (t,x,y,r) or (x,y,z,w).
\section{因子変形を持ち上げる二つの機構}\label{sec:moving-examples}
四次元で$\alpha=\iota_P\mu$、$q(\alpha)=\alpha^2/(2\mu)$とし、
$d_f\alpha=f\,d\alpha-df\wedge\alpha$と書く。
局所体積の下で$\alpha_\epsilon=\alpha+\epsilon a+\epsilon^2b+\cdots$、
$f_\epsilon=f+\epsilon\varphi+\epsilon^2\psi+\cdots$なら、必要十分な二次までの式は
\begin{align}
 \alpha\wedge a/\mu&=\varphi,&
 d_fa&=d\varphi\wedge\alpha-\varphi d\alpha,\label{eq:example-first}\\
 \alpha\wedge b/\mu&=\psi-q(a),&
 d_fb&=d\psi\wedge\alpha-\psi d\alpha
                  +d\varphi\wedge a-\varphi da.\label{eq:example-second}
\end{align}
これは$q(\alpha_\epsilon)=f_\epsilon$と$d_{f_\epsilon}\alpha_\epsilon=0$の展開である。
以下の二例はこれを全次数で満たす実族であり、任意の因子方向の無障害性は主張しない。

\subsection{TRIPLEからGa9：二次補正が強制される理由}
座標$(t,x,y,r)$、$\mu=dt\wedge dx\wedge dy\wedge dr$とし、中心を
\[
\begin{gathered}
 P_0^{tx}=x/3+y,\quad P_0^{ty}=-xy+y/6-r/3,\quad
 P_0^{tr}=-3y^2+xr-5y/4+r/2,\\
 P_0^{xy}=1,\quad P_0^{xr}=-3x+3/2,\quad P_0^{yr}=3x^2+3y
\end{gathered}
\]
とする。$H=x^3+3xy-3y/2$、$k=xy/3+y^2/2-x^4/12$、
$K=k+x(r+H)/3$と置けば$f_0=r+H$、$Z_0=-\partial_t$である。
直積括弧$\{x,y\}=1,\{T,s\}=s$から
\[
 r=s+eT-H,\qquad
 t=(1-ex/3)T-xs/3-k+M,\qquad M=A(e)x^2y+B(e)y^2
\]
を取る。この写像のJacobianは全eで一で、逆は
$T=t+x(r+H)/3+k-M$、$s=r+H-eT$。
従って$1-ex/3$を反転する必要はない。
ただし$M=0$のままでは$P^{tr}$に三次以上の項が現れ、二次Poisson空間を出る。
実際
\[
 P^{tr}=P_0^{tr}-et-eK+eM-\{M,H\},\quad
 \{M,H\}=-3Ax^4+(3A-6B)x^2y-3Axy-6By^2.
\]
$x^4,x^2y$を消す条件はそれぞれ
$-e/4+3A=0$、$-e+eA-3A+6B=0$であり、このansatz内で一意に
\[
 A=e/12,\qquad B=5e/24-e^2/72
\]
を強制する。従って補正の理由は因子の既約化そのものではなく、
Poissonである押出しを\emph{係数次数二以下に保つ}ことである。

$P_e=P_0+eP_1+e^2P_2+e^3P_3$の非零補正成分は次の通りで、未表示成分は零である。
\[
\begin{array}{c|ccc}
 &tx&ty&tr\\\hline
 P_1&-5y/12-x^2/12&xy/6&-t-xr/3+3y^2/4+5xy/12\\
 P_2&y/36&0&y^2/8\\
 P_3&0&0&-y^2/72
\end{array}
\]
Pfaffianとmodular場は全eで
\begin{align*}
 f_e&=f_0+e\varphi+e^2\psi+e^3\chi,\\
 \varphi&=-t-K,\qquad
 \psi=x^2y/12+5y^2/24,\qquad\chi=-y^2/72,\\
 Z_e&=Z_0+eZ_1,\qquad Z_1=(x/3)\partial_t-\partial_r
\end{align*}
である。$a=\iota_{P_1}\mu,b=\iota_{P_2}\mu$とすると、
命題\ref{prop:modular-primitive}の一次・二次式は
\[
 d\varphi+\iota_{Z_0}a+\iota_{Z_1}\alpha_0=0,\qquad
 d\psi+\iota_{Z_0}b+\iota_{Z_1}a=0.
\]
これらとPfaffian展開は式\eqref{eq:example-first}--\eqref{eq:example-second}を満たす。
一方$P_1$は$\partial_t$で割れるので$q(a)=0$だが$\psi\ne0$。
従って\emph{この指定因子族と一次持上げ}に対し$b=0$は不可能であり、
$\alpha_0\wedge b/\mu=\psi$が必要である。
実際$P_0+eP_1$だけのJacobiの$(t,x,r)$成分には$e^2(x^2+5y)/12$が残る。
これは全二次障害空間が非零という結論ではなく、表示した持上げの補正項を捨てられないという結論である。

$f_e$のt係数は$-e$、四次係数は$-ex^4/4$である。
従って$e\ne0$で$f_e$はtについて非零定数係数の一次式で絶対既約、$\deg f_e=4$。
$e=0$では$\deg f_0=3$なので、actual因子は
$2H+C_3$から単純$H+F_4$へ変わる。
$P_e^{xy}=1$により全eで射影tensorは非零で、$f_e$も全eで非零である。
この族がGa9成分へ入ることには、別の実アフィン変更
\[
 (X,Y,Z,w)=(t-r/e,ex/3,-r/e,3y/e),\qquad \det=1/e
\]
と既存の\emph{全十源}への係数同定を使う。
元証明\path{CONSTANT_MODULAR_TRIPLE_WHOLE_GA9_CONSTANT_ACTUAL_FAMILY.md}第5節の
$r_{\mathrm{src}}=81/e^4\ne0$、$F_{2,\mathrm{src}}=e^2(e-3)(e-12)/648$を用い、
$e=3,12$の零補集合も保持する。
非線形の積模型だけからPGL帰属を推論していない。
残る全forced sourceへの拡張も、同ノート第1節の実normalizerによる全源から一軌道への還元を入力とする。

\subsection{二重Hを分裂する可換場の族：高次補正が不要な場合}
座標$(x,y,z,w)$で、互いに可換な線形場
\[
 X=x\partial_x+y\partial_y-2z\partial_z,\quad
 Y=x\partial_x-y\partial_y,\quad D=z\partial_z
\]
を取り
\[
 P_e=X\wedge Y+(1+ew)D\wedge\partial_w
\]
と置く。$X,Y$は$D,\partial_w,w$を保存するから全eでJacobiは零である。
全六係数とPfaffianは
\[
 P_e=(-2xy,2xz,0,-2yz,0,z(1+ew)),\qquad
 f_e=-2xyz(1+ew).
\]
$a=\iota_{wD\wedge\partial_w}\mu$は分解可能なので$q(a)=0$であり、
$[P_0,wD\wedge\partial_w]=0$と$[wD\wedge\partial_w,wD\wedge\partial_w]=0$も成立する。
従って式\eqref{eq:example-first}の$\varphi=-2xyzw$に対して
$b=\psi=0$が可能で、以後全ての補正も零でよい。
actual五次式は$-2hxyz(h+ew)$で、中心は$2H+H_x+H_y+H_z$、
$e\ne0$では五つの独立超平面に分裂する。
このとき$w'=1+ew$は真の可逆アフィン変更で、$\partial_w=e\partial_{w'}$となり、
$P_e$は非退化logcanonical型へ移る。従ってT01へ入り、閉性から中心もT01に入る。
これは可換場の全源族を使う
\path{MULTIPLE_HYPERPLANE_MODULAR_SUPPORT.md}第3節の機構の明示例である。
TRIPLEの次数補正とは異なり、可換性と補正方向の分解可能性が全高次障害を消す。
一般の因子分裂にもこの性質があるとは主張しない。


\subsection{閉包とrank二境界を正しく扱う}
一般階数四の開集合を$U_4$、そのPoisson空間内の閉包を$Y$（被約構造）とする。
$q$の全係数が生成するイデアル層を$J$と書く。

\begin{proposition}[Pfaffianグラフの閉包]\label{prop:graph}
$\delta$のグラフ閉包$\Gamma\subset Y\times\PP(\mathbb B)$は
$\operatorname{Proj}_Y\bigoplus_{n\ge0}J^n$と同一視できる。
第一射影の像は$Y$全体であり、一般rank二の独立成分を復元するものではない。
\end{proposition}
\begin{proof}
局所的に$q$の係数を$q_0,\ldots,q_N$と書く。
Rees代数への準同型$\cO_Y[T_0,\ldots,T_N]\to\cO_Y[t]$、$T_i\mapsto q_it$が
射影閉部分を定め、$J\ne0$の開でグラフと一致する。
$Y$は被約で$U_4$が稠密なので、この代数の正次数部分に$U_4$を避ける余分な既約成分はない。
したがってそのProjはグラフ閉包である。これはblowupのRees表示でもある\cite{StacksBlowup}。
射影性から第一射影の像は閉で、$U_4$を含むため$Y$である。
\end{proof}

生のincidence方程式で$q=0$も許すと、このグラフ閉包とは異なる集合を得る可能性がある。
グラフ閉包はrank四開から取ることが本質であり、計算ではその開に関する飽和が必要である。
ただし\emph{底の特別な因子層}にだけあるrank四成分を、一般$F$へ優越しないという理由で捨ててはならない。
$tx=0$で与えられる族では、$t\ne0$のkernelは零だが$t=0$に一本の余分な成分がある。
定数階数層の補集合も研究対象であることを示す最小の例である。
全Reesイデアルを巨大なGr\"obner計算で求めることを、この構成の実行要件とはしない。

\subsection{標識因子と一般ファイバーの重複を管理する}
\begin{proposition}\label{prop:factor-mark}
正整数対$(d_i,m_i)$、$\sum_i d_im_i=5$を固定する。
積写像
\[
 \prod_i\PP H^0(\PP^4,\cO(d_i))\longrightarrow\PP(\mathbb B),
 \qquad([f_i])\longmapsto[\prod_i f_i^{m_i}]
\]
はその閉像への有限写像である。
\end{proposition}
\begin{proof}
射影空間の積からの射なので固有である。非零多項式の一意分解により、
指定された積に対する因子の割当ては有限個であり、各因子のスカラーは射影化で除かれる。
従って準有限でもあり有限となる。
\end{proof}

因子を標識したパラメータ空間は境界での衝突・分裂も保持できるが、標識を忘れた重複を数えてはならない。
さらに、既約な底$S$上の復元族$T\to S$のうち$S$を優越する既約成分は、
$\CC(S)$上の一般ファイバーの既約成分と対応する。
これは局所環を$\cO_S\setminus\{0\}$で局所化した素イデアルの対応から従う。
幾何学的一般ファイバーを使うときはGalois軌道をまとめる必要がある\cite{StacksFibres}。
例えば$x^2=t$は全空間が既約だが、幾何学的一般ファイバーは二点である。
特別な底に支持される成分は別に残し、有限個の閉既約像の全点被覆が得られた後、
\emph{その中の極大な相異なる閉像}を成分として数える。
これは全軌道分類や全成分対の交叉計算を要求しない。

\section{階数二はfoliationと零因子から復元する}
\begin{proposition}[階数二の構造的入口]\label{prop:ranktwo}
$\PP^4$上の非零で一般階数二のPoisson構造は、飽和した階数二の可積分部分層
$\mathcal F\subset T_{\PP^4}$と、非零切断
$s\in H^0(\PP^4,\det\mathcal F)$の組として一意に記述される。
ここで$\det\mathcal F=(\wedge^2\mathcal F)^{**}$、
$\det\mathcal F=\cO(e)$なら$0\le e\le2$である。
$s$の零因子の次数は$e$、foliationの次数を$d=2-e$とすれば$d+e=2$となる。
射影化されたPoisson構造では$s$の非零スカラー倍を同一視する。
\end{proposition}
\begin{proof}
$\pi^\sharp$の像を$T$内で飽和して$\mathcal F$を得る。
一般点で$\pi$は分解可能であり、Jacobi恒等式から像分布は可積分である。
Lie括弧を$T/\mathcal F$へ写す障害は一般点で零で、商はtorsion-freeなので全域で零となる。
飽和部分層$\mathcal F$は反射的であり、余次元一では局所自由な部分束となる。
その外積は余次元一で原始的なので、$\pi$は$\det\mathcal F$の正則切断$s$を定める。
余次元二以上へは反射性によって延長し、零因子も一致する。

逆に飽和した可積分$\mathcal F$について、余次元一の外積写像
$\det\mathcal F\to\wedge^2T$は反射性により全域へ延長する。
非零$s$の像$\pi$は一般階数二であり、局所自由な開で$\pi=aX\wedge Y$と書けば
$[X,Y]\in\mathcal F$から$[\pi,\pi]=0$が従う。正則テンソルの等式として全域へ延長する。
両構成は一般点で互いに逆であり、飽和性・反射性から全域でも逆である。

$s\ne0$から$e\ge0$。外積写像は
$H^0(\wedge^2T(-e))=H^0(\Omega^2(5-e))$の非零元を与える。
射影空間上の二形式の同次水平表示では、$H^0(\Omega^2(k))=0$が$k\le2$で成り立つ。
実際、$k<2$なら係数次数$k-2$が負であり、$k=2$なら定数二形式のEuler収縮が零になるのは零形式だけである。
従って$5-e\ge3$、すなわち$e\le2$である。
\end{proof}

この命題はrank二の全入口を$d=0,1,2$のfoliationと零因子へ分けるが、
各foliation族の認識や全Poisson空間内での成分性は証明していない。
rank二部分空間の成分と、全空間の一般rank二成分を区別する。
一般rank四の閉包に入るrank二族は命題\ref{prop:graph}の側で処理し、
純粋なrank二成分は上のfoliationの支持定理から接続する。

\section{構造原理が説明する代表的な現象}
\subsection{五つの座標超平面と25次元族}
$D_0=V(x_0x_1x_2x_3x_4)$とし、$x_0\ne0$上で$y_i=x_i/x_0$と置く。
任意の非退化な交代行列$A=(a_{ij})_{1\le i,j\le4}$に対し
\[
 \pi_A=\sum_{i<j}a_{ij}(y_i\partial_{y_i})\wedge(y_j\partial_{y_j})
\]
は同じ退化因子$D_0$を持つPoisson構造である。
対数接束$T_{\PP^4}(-\log D_0)$は四つのトーラス場で自明化されるので、逆にこの因子を持つ全構造はこの形である。
$A$の自由度は6、射影化後は5であり、$D$だけでは$\pi$は決まらない。
$D_0$の射影自己同型群の単位成分は四次元トーラスで、すべての$\pi_A$を保存する。
従って軌道次元は$24-4=20$であり、因子を動かした族の次元は$20+5=25$となる。
この次元計算だけから全Poisson空間の成分性を結論するわけではない。

正規交叉の退化因子を持つ射影空間の構造がトーリックであることは\cite[Proposition 5.16]{Pym}の既知結果である。
SNC点から出るより多くの成分をsmoothing diagramと共鳴で扱うのが\cite{MPS}であり、
同論文の40成分という結果を全Poisson空間の尽くしと読み替えない。

\subsection{二因子型と三因子型}
$D=H+Q_4$や$D=2H+C_3$は、重複度が違っても被約因子数は$r=2$である。
定理\ref{thm:abelian}は一つの非零な線形対称性を与える。
そのJordan分解、半単純部分の重み、冪零部分の旗をまず調べることが、個別のcofactor計算の共通の理由となる。
一方$D=H_1+H_2+C_3$は$r=3$であり、二次元の可換対称性が強制される。

軌道23次元の二因子型と、軌道22次元に一つの本質パラメータが加わる三因子型は、
ともに23次元族になり得る。この説明は147成分台帳の個々の成分を新しく証明するものではないが、
似た次元が反復する理由を整理する。
追加の対称性があるか、残りのパラメータが真に非同型を与えるかは、それぞれ確認する必要がある。


\begin{proposition}[modular場との対応]\label{prop:modular}
アフィン四空間上の一般に非退化なPoisson双ベクトルを$P$とし、
標準体積形式$\mu$に関して$P^2/2=F\mu^{-1}$と書く。
$Z^j=\sum_i\partial_iP^{ij}$で定めたmodular場は
\[
 Z=\frac{P^\sharp(dF)}{F}.
\]
従って$F=c\prod_i g_i^{m_i}$なら
$Z=\sum_i m_iP^\sharp(dg_i)/g_i$である。
特に$D=m_0H+m_1D_1$で$H$を無限遠に取る場合、$Z$は
定理\ref{thm:abelian}の一次元対称性代数の生成元の非零定数倍である。
\end{proposition}
\begin{proof}
$H_f=P^\sharp(df)$とすると、成分を微分して
$\diver_\mu H_f=-Z(f)$を得る。
一方、シンプレクティック体積形式$\mu/F$はHamiltonian場で保存されるので
$\diver_\mu H_f=H_f(F)/F=-P^\sharp(dF)(f)/F$である。
両式を比べれば最初の恒等式となる。因数分解については対数微分を取ればよい。
最後に$D_1$の次数を$d$とし、$H=V(h)$とする。
$W_D$の生成元に対応する形式は$dh/h-d^{-1}df_1/f_1$である。
$h=1$上でこれは$-d^{-1}dg_1/g_1$となり、
$Z=m_1P^\sharp(dg_1)/g_1$と比例する。
非零性は定理\ref{thm:abelian}の単射性から従う。
\end{proof}

\subsection{重複二次因子と共形シンプレクティック線形代数}
実因子$D=H+2Q$で、不変超平面$H$を無限遠とする二次Poisson代表が得られたとする。
アフィン四空間上で逆形式を$\omega=\alpha/Q^2$と書くと、定理\ref{thm:fibre}の条件は
\[
 Qd\alpha-2dQ\wedge\alpha=0,\qquad
 \alpha^2/2=kQ^2\mu,\quad k\ne0.                    \tag{7.1}
\]
もし全閉形式が$\alpha=Qd\beta-dQ\wedge\beta$、$\deg\beta\le1$と表示できるなら、
$\omega=d(\beta/Q)$である。ここで$\Omega=d\beta$が非退化なら
$\iota_Z\Omega=\beta$によってアフィン線形場$Z$を定める。
$\Omega^2/2=A\mu$、$A\ne0$と置けば
\[
 \mathcal L_Z\Omega=\Omega,\qquad
 \alpha^2/2=AQ\bigl(Q-Z(Q)\bigr)\mu.
\]
従って(7.1)は
\[
 \mathcal L_Z\Omega=\Omega,\qquad
 Z(Q)=cQ,\quad c=1-k/A\ne1                          \tag{7.2}
\]
へ還元される。
\begin{proof}[最後の恒等式の証明]
$\alpha=Q\Omega-dQ\wedge\beta$を二乗する。
$(dQ\wedge\beta)^2=0$であり、
$\Omega\wedge dQ\wedge\beta=dQ\wedge\iota_Z(\Omega^2/2)=A Z(Q)\mu$となる。
また$\mathcal L_Z\Omega=d\beta=\Omega$である。
\end{proof}

これは90係数の方程式を、二次式を半不変量に持つ共形シンプレクティック場の分類へ置き換える仕組みである。
ただし「全閉形式が上の$\beta$で表示できる」ことと、$d\beta$が退化する場合の処理は別の入力である。
現行の\texttt{DOUBLE\_QUADRIC\_SUPPORT.md}はこの仕組みを用いているが、
本稿ではそこでの有限行列階数計算を概念的に置き換え終えたとは主張しない。
この入力の概念化は有用だが、大域分類の最優先課題とはしない。
とくに射影階数三の横断配置では低次数primitiveの像に対する余剰は一、接配置では六次元であり、
無条件の一様primitive表示は偽である。前者には非零de Rham類、後者にはprimitiveの次数制約が関係する。
これらの次元は既存の有理数上の線形計算に依存し、本稿では再計算していない。
原理の適用で必要になった範囲だけ、次数・極位数付きの複体として説明する。

\section{成分分類の再構成に必要な三つの段階}
\subsection{構造：どの幾何が可能か}
まず一般階数四と二を分ける。階数四について$D$の次数分割と重複度を保持し、
$\mathfrak a_\pi$の表現、$D_i$上の葉、必要なら正規化と導手を調べる。
$r=1$では本稿の可換代数は零だが、既約五次と$5H$を区別する。
前者には正規・非正規の幾何、後者には系\ref{cor:unimodular}のunimodular二次模型を用いる。
$r\ge2$では強制される対称性を先に用い、全可換部分代数の無制限な分類ではなく、
次数五の因子と二次Poisson係数が許す表現だけを扱う。

\subsection{復元：同じ幾何の上に何が残るか}
固定$D$について$V_F$を決定し、その中のPfaffian二次条件を解く。
この解には$D$を固定したまま変わる自由度も残るので、因子だけの分類をPoisson構造の分類と同一視しない。
閉形式のコホモロジー類だけでも十分とは限らず、極の主要部と大域への延長条件を保持する。

\subsection{変形：どれが別成分か}
上の型が有限の閉既約族で全点を覆うことを示し、その後に重複・包含を整理する。
一般元の因数分解型やJordan型が違っても、一方が他方の境界に入る可能性はある。
退化族、閉incidence条件、変形の障害、葉の次数や対数留数を用いて包含・非包含を証明する。
同じ型の境界を新成分と数えない。

主稿の推奨順序は「全入口と標識」「共通の復元・変形原理」「超平面型」
「二次＋三次型」「既約五次の認識」「全支持と極大閉像」「rank二の支持」「全成分表」である。
134方程式や$S_{61}$はこの構造を実装・検証する共通座標模型として保持し、
一般読者が最初に全座標を読まなければ主定理の意味が分からない構成を避ける。
147という個数は構造と包含関係を証明した\emph{結果}として現れるべきであり、整理の前提にしない。

\section{主張の状態と先行研究}
対称性・固定因子復元に加え、定理\ref{thm:three-entrances}、\ref{thm:marked-affine}、
\ref{thm:modular-reconstruct}、命題\ref{prop:resonance}、\ref{prop:lifting}、\ref{prop:graph}、\ref{prop:ranktwo}には本稿で証明を与えた。
これらは既存のPoisson幾何の基本的道具を組み合わせるものであり、新規性を主張しない。
特に対数Hamiltonian場は\cite[Example 4.14]{MPS}にも現れる。
本稿の非被約因子も許す大域的$r-1$次元表示がどこまで既述かについて、文献の完全な照合は未完了である。

Pymの\cite[\S5.1]{Pym}では退化因子のPoisson性と逆対数形式、
\cite[Propositions 5.14, 5.16]{Pym}ではSNCの場合のChern指標制約を説明する。
MPSの\cite[\S3.1, Theorem 1.4]{MPS}はSNCに関連する変形とsmoothing diagramを扱う。
いずれも本稿の二つの原理だけで147成分の尽くしが従うという主張を含めて引用しない。

MPSの定理はSNC locusに交わる40成分を扱い、全成分がトーリック型へ退化するという主張ではない。
階数二の入力についても、Corr\^ea--Munizの本文Theorem 4.1\cite{CM}が安定分岐に与えるのは
代数的可積分性とuniruledな葉であり、これだけから固定されたJacobian族への全点帰属を引用しない。

共通ファイルを主分類稿にも読み込み、単純Hと重複Hの標識・modular・共鳴の入口の重複を置換した。
modular原始関数の全恒等式、TRIPLEの次数補正と指定因子族の二次持上げ、可換場による分裂を確認した。
二次＋三次の二つのcofactor入口、平面型の三つの全source復元、
二因子と既約因子の磁気shearの次元証明も共有TeXへ統合した。
各源の帰属・成分性・非包含の旧証明は保持し、全147成分を独立再検証したとはしない。
検算はそれらの構造定理に必要な有限恒等式へ位置づけ、研究の主目標そのものとはしない。

\begin{thebibliography}{9}
\bibitem{Pym} B. Pym, \emph{Constructions and classifications of projective Poisson varieties},
Letters in Mathematical Physics \textbf{108} (2018), 573--632.
DOI: 10.1007/s11005-017-0984-5; arXiv:1701.08852.
\url{https://pmc.ncbi.nlm.nih.gov/articles/PMC5818673/}.
\bibitem{MPS} M. Matviichuk, B. Pym and T. Schedler,
\emph{A local Torelli theorem for log symplectic manifolds},
arXiv:2010.08692v2 (15 July 2021).
\url{https://arxiv.org/html/2010.08692v2}.
\bibitem{CM} M. Corr\^ea and A. Muniz,
\emph{Holomorphic foliations of degree two and arbitrary dimension},
arXiv:2207.12880v5, Theorem 4.1.
\url{https://arxiv.org/html/2207.12880v5}.
\bibitem{StacksBlowup} The Stacks Project, Tag 01OF, Blowing up.
\url{https://stacks.math.columbia.edu/tag/01OF}.
\bibitem{StacksFibres} The Stacks Project, Tag 0553, Irreducible components of fibres.
\url{https://stacks.math.columbia.edu/tag/0553}.
\end{thebibliography}
\end{document}
